% !TITLE 观沧海
% !AUTHOR 曹操
% !DYNASTY 魏晋
% !GENRE 四言古诗
% !ORDER 3

\begin{poem}
东临碣石，以观沧海。
水何澹澹，山岛竦峙。
树木丛生，百草丰茂。
秋风萧瑟，洪波涌起。
日月之行，若出其中；
星汉灿烂，若出其里。
幸甚至哉，歌以咏志。
\end{poem}

\begin{appreciation}
这是中国文学史上第一首完整的山水诗，也是曹操最著名的诗作，写于他北征乌桓得胜回师途中。

建安十二年（207年），曹操率军北征乌桓，消灭袁绍残余势力，统一了北方。回师途中，他登上碣石山，面对渤海，写下了这首气吞山河的诗。

东临碣石，以观沧海——向东登上碣石山，来观赏大海。开篇直接交代地点和目的，没有铺垫，没有修饰，像一个军事统帅简洁明快的命令。这就是曹操的风格：不拖泥带水，直奔主题。

水何澹澹，山岛竦峙——海水多么浩荡，山岛高高耸立。澹澹（dàn dàn）是水波摇动的样子，竦峙（sǒng zhì）是高耸挺立的样子。这十个字勾勒出一幅壮阔的海景图：无边无际的海水，突兀而起的岛屿，动静相映，远近交错。

树木丛生，百草丰茂。秋风萧瑟，洪波涌起——岛上树木茂密，百草繁盛。秋风吹过，波涛汹涌。这里有一个微妙的转折：前两句写静态的景物，后两句写动态的秋风洪波。秋天本是万物凋零的季节，但在曹操眼中，秋天的沧海却充满了力量。洪波涌起四个字，写出了大海在秋风中的磅礴气势，也暗示了诗人内心的激荡。

日月之行，若出其中；星汉灿烂，若出其里——太阳和月亮的运行，仿佛出自大海之中；灿烂的银河，仿佛出自大海之内。这是全诗最宏大的想象。曹操不是站在海边看海，而是把大海看作宇宙的中心。日月星辰都在大海中起落，这是何等的气魄！这种以宇宙为背景的视角，是中国诗歌中少见的。

最后两句幸甚至哉，歌以咏志是乐府诗的套语，每首乐府结尾都可以加，与内容无关。真正的诗眼在前面的十句。读这首诗，你能感受到一个政治家兼诗人的胸襟——不是小儿女的情长，不是文人的酸腐，而是一个统一天下的英雄面对自然时发出的惊叹。这种气概，后世称为建安风骨。
\end{appreciation}
